\section{Discussion}
The results indicate that the proposed system is best understood as a first-pass filter, not as a complete hiring intelligence platform. It can rapidly identify candidates who clearly lack explicit qualifications while preserving an explanation trail. For high-volume job postings, this staged approach is operationally useful: a lightweight lexical system can perform the initial pass quickly, while human review or semantic matching can be reserved for a smaller shortlist.

The explanation mechanism is deliberately simple. Since every score component is computed from an explicit field, the system directly states which requirements were matched and which were missing. This style of explanation is suited to HR users who need actionable justification rather than model-internals visualization. It also helps maintenance because extraction errors can be traced to the skill dictionary, parser output, or normalization rules.

The default architecture avoids GPU dependencies and external API calls, making it practical for classroom demonstration, small HR teams, or early-stage organizations. The OCR fallback contributes to fairness by reducing format bias, but fairness in recruitment is broader than file handling \cite{b4}. The system should therefore be used to prioritize review, not to reject candidates without human oversight. The optional hybrid context mode demonstrates that semantic embeddings can be added while preserving deterministic hard-skill checks, although this increases runtime and deployment requirements.

\section{Limitations}
The most important limitation is the lexical gap. The system may fail to reward semantically equivalent phrasing if the exact skill token is absent. For example, ``RESTful services'' may not match a strict requirement for ``REST'' depending on ontology design. Similarly, soft-skill evidence may be expressed through narrative rather than direct words such as ``leadership''.

The skill ontology is curated and therefore requires maintenance. New technologies, abbreviations, and synonyms must be added to keep the system current. The OCR fallback improves robustness, but OCR can still fail on complex layouts, handwritten content, low-resolution scans, or heavily graphical resumes. The benchmark dataset is also limited and should be expanded with real-world resumes across industries and formats.

Finally, the system measures document-role alignment, not candidate quality in a complete sense. It cannot evaluate interview performance, portfolio depth, communication ability, or cultural fit. Its purpose is to support recruiters in prioritizing review, not to replace human decision-making.

\subsection{Validity Considerations}
The benchmark suite is focused and controlled, but it is not a large industrial dataset. Results may differ on highly visual templates, multilingual documents, or specialized industries. The system also depends on ontology coverage; missing aliases such as ``Postgres'' for ``PostgreSQL'' can reduce recall. OCR quality depends on image resolution, layout complexity, and scan noise. Finally, timing results should be interpreted as prototype measurements rather than universal constants, although the relative efficiency advantage of TF-IDF over embedding models remains technically expected.

\section{Future Work}
Future work can improve the framework in four directions. First, a controlled synonym expansion layer can reduce lexical brittleness while preserving deterministic explanations. Second, the hybrid context mode can be calibrated with larger embedding models, caching strategies, and threshold studies to determine when semantic comparison improves ranking enough to justify additional cost. Third, the skill ontology can be moved from static dictionaries to an updateable skill graph with aliases and related technologies. Fourth, fairness evaluation can be expanded using larger datasets with format diversity, role diversity, and adversarial resume patterns.

An additional future direction is role-specific calibration. Different job families require different scoring behavior. Engineering roles may emphasize technical skills, while management roles may require more nuanced weighting for communication, leadership, and years of experience. The current weight sliders provide an initial mechanism, but future systems can learn recommended weight profiles from recruiter feedback.

\section{Conclusion}
This paper presented an explainable and cost-efficient AI resume screening framework for first-pass recruitment filtering. The system combines OCR-aware parsing, rule-guided skill extraction, switchable contextual similarity, bounded weighted scoring, and plain-language explanations. Experimental benchmarks show that the framework handles multiple file formats, resists simple keyword-stuffing behavior, preserves mathematical score boundaries, and processes resume batches with very low time and memory requirements in its default configuration. The approach does not eliminate the semantic advantages of transformer models; instead, it exposes a controlled hybrid option for cases where semantic context is more important than minimum latency. The proposed framework is therefore suitable as an auditable screening layer that narrows applicant pools before human review or deeper semantic analysis.
